\documentclass[12pt, a4paper]{article}

% --- PREÁMBULO ---
% PAQUETES BÁSICOS
\usepackage[utf8]{inputenc}
\usepackage[spanish]{babel}
\usepackage{amsmath}
\usepackage{geometry}
\usepackage{enumitem}
\usepackage{fancyhdr}
\usepackage{graphicx}

% PAQUETES PARA TABLAS MEJORADAS
\usepackage{booktabs} 
\usepackage{array}    

% CONFIGURACIÓN DE LA PÁGINA
\geometry{
    a4paper,
    total={170mm,257mm},
    left=20mm,
    top=20mm,
    right=20mm,
    bottom=20mm,
    headheight=15pt
}

% CONFIGURACIÓN DEL ENCABEZADO
\pagestyle{fancy}
\fancyhf{}
\fancyhead[L]{126}
\fancyhead[R]{Prácticas de electricidad}
\renewcommand{\headrulewidth}{0.4pt}

% --- CUERPO DEL DOCUMENTO ---
\begin{document}

% --- TABLA 20-1 ---
\begin{table}[h!]
    \centering
    \caption{TABLA 20-1}
    \label{tab:20-1}
    \begin{tabular}{lcccccccc}
        \toprule
        & $R_1$ & $R_2$ & $R_3$ & $R_4$ & $R_5$ & $R_6$ & $R_7$ & $R_8$ \\
        \midrule
        \begin{tabular}{@{}l@{}}Valor nominal \\ ohmios\end{tabular} & 330 & 470 & 1.200 & 2.200 & 3.300 & 4.700 & 10.000 & 5.600 \\
        \bottomrule
    \end{tabular}
\end{table}

% --- TABLA 20-2 ---
\begin{table}[h!]
    \centering
    \caption{TABLA 20-2}
    \label{tab:20-2}
    \begin{tabular}{lccccccccccc}
        \toprule
        & $I_T$ en A & $I_2$ & $I_3$ & $I_T$ en B & $I_T$ en C & $I_5$ & $I_6$ & $I_7$ & $I_T$ en D & $I_2+I_3$ & \begin{tabular}{@{}c@{}}$I_5+I_6$\\$+I_7$\end{tabular} \\
        \midrule
        mA & & & & & & & & & & & \\
        \bottomrule
    \end{tabular}
\end{table}

% --- TABLA 20-3 ---
\begin{table}[h!]
    \centering
    \caption{TABLA 20-3}
    \label{tab:20-3}
    \begin{tabular}{lcccccc}
        \toprule
        & $I_T$ en A & $I_1$ & $I_2$ & $I_3$ & $I_T$ en B & $I_1+I_2+I_3$ \\
        \midrule
        mA & & & & & & \\
        \bottomrule
    \end{tabular}
\end{table}

% --- TABLA 20-4 ---
\begin{table}[h!]
    \centering
    \caption{TABLA 20-4}
    \label{tab:20-4}
    \begin{tabular}{lccccc}
        \toprule
        & $I_1$ & $I_2$ & $I_3$ & $I_T$ & $I_1:I_2:I_3$ \\
        \midrule
        mA & & & & & \\
        \bottomrule
    \end{tabular}
\end{table}

% --- PREGUNTAS ---
\section*{PREGUNTAS}
\begin{enumerate}[label=\arabic*., start=2]
    \item ¿Cómo es la suma de $I_5$, $I_6$ e $I_7$ del apartado 1 comparada con $I_T$ en C? ¿$I_T$ en D? Explicar las diferencias.
    \item Comentar los resultados de las propias mediciones del apartado 2. Explicar cualquier resultado inesperado.
    \item Enunciar la ley de corriente de Kirchhoff y escribir su fórmula.
    \item ¿Qué información es necesaria para calcular $I_2$ e $I_3$ en la figura 20-4?
    \item Discutir los resultados del problema de diseño (apartado 3).
    \item ¿Por qué no es posible diseñar un circuito que dé una relación exacta de corriente 1:2:3?
\end{enumerate}

% --- RESPUESTAS AL AUTOEXAMEN ---
\subsection*{Respuestas a las preguntas de autoexamen}
\begin{enumerate}[label=\arabic*.]
    \item 0,2.
    \item 1,25.
    \item $I_1$ es +; $I_2$ es +; $I_3$ es +; $I_T$ es –.
    \item $I_1 + I_2 - I_3 - I_4 - I_5 = 0$.
    \item 2.
\end{enumerate}

\end{document}
